\documentclass[../../../topic_linear-algebra]{subfiles}

\usepackage{xr-hyper}
\externaldocument{../../../.tex_intermediates/topic_linear-algebra}

\begin{document}

\sectionline
\section{部分空間の和}
\marginnote{\refbookS p20〜21 \\ \refbookR p47〜48 \\ \refbookC p22〜23 \\ \refbookF p231〜232 \\ \refbookG p6}

ある線形空間を、複数の線形空間の「和」に分解して考えることが、この章の一つめの主題である。

\br

まずは素朴に、2つの部分空間$U, W$の\keyword{和}$U + W$を考える。

これは、$U$の元と$W$の元を足し合わせた全ての元からなる集合である。

\begin{theorem*}{部分空間の和の部分空間性}
  $U,W$が$V$の部分空間であるとき、\keywordJE{和}{sum}$U + W$は$V$の部分空間である。
  \begin{equation*}
    U + W = \{ \vb*{u} + \vb*{w} \mid \vb*{u} \in U, \vb*{w} \in W \}
  \end{equation*}
\end{theorem*}

\begin{proof}
  \begin{subpattern}{\bfseries 和について}
    $\vb*{a}_1, \vb*{a}_2 \in U, \, \vb*{b}_1, \vb*{b}_2 \in W$とする。

    $U$と$W$は部分空間なので、部分空間の定義より、
    \begin{equation*}
      \vb*{a}_1 + \vb*{a}_2 \in U, \quad \vb*{b}_1 + \vb*{b}_2 \in W
    \end{equation*}

    一方、和空間の定義より、$\vb*{a}_1 + \vb*{b}_1,\, \vb*{a}_2 + \vb*{b}_2$はそれぞれ$U+W$の元である。

    これらの元の和をとったときに、その和も$U + W$に属していれば、和空間は和について閉じているといえる。
    \begin{equation*}
      (\vb*{a}_1 + \vb*{b}_1) + (\vb*{a}_2 + \vb*{b}_2) = (\vb*{a}_1 + \vb*{a}_2) + (\vb*{b}_1 + \vb*{b}_2)
                                                        \in U + W
    \end{equation*}
    上式で、和空間は和について閉じていることが示された。 $\qed$
  \end{subpattern}

  \begin{subpattern}{\bfseries スカラー倍について}
    $\vb*{a} \in U, \, \vb*{b} \in W$と$c \in K$をとる。

    $U$と$W$は部分空間なので、部分空間の定義より、
    \begin{equation*}
      c \vb*{a} \in U, \quad
      c \vb*{b} \in W
    \end{equation*}

    一方、和空間の定義より、$\vb*{a} + \vb*{b}$は$U + W$の元である。

    この元をスカラー倍したときに、そのスカラー倍も$U + W$に属していれば、和空間はスカラー倍について閉じているといえる。
    \begin{equation*}
      c(\vb*{a} + \vb*{b}) = c \vb*{a} + c \vb*{b}
                           \in U + W
    \end{equation*}
    上式で、和空間はスカラー倍について閉じていることが示された。 $\qed$
  \end{subpattern}
  
  \begin{subpattern}{\bfseries 零元の存在}
    $U$と$W$は部分空間なので、零元が存在する。
    \begin{equation*}
      \vb*{o} \in U, \quad \vb*{o} \in W
    \end{equation*}
    よって、
    \begin{equation*}
      \vb*{o} = \vb*{o} + \vb*{o} \in U + W
    \end{equation*}
    より、和空間にも零元が存在することが示された。 $\qed$
  \end{subpattern}
\end{proof}

\subsection{複数の部分空間の場合}

より一般に、複数の部分空間$W_1, \ldots, W_n$の\keyword{和}も部分空間となる。

\begin{theorem*}{部分空間の和の部分空間性（一般の場合）}
  $W_1, \ldots, W_n$が$V$の部分空間であるとき、\keywordJE{和}{sum}$W_1 + \ldots + W_n$は$V$の部分空間である。
  \begin{equation*}
    W_1 + \ldots + W_n = \{ \vb*{w}_1 + \ldots + \vb*{w}_n \mid \vb*{w}_i \in W_i \}
  \end{equation*}
\end{theorem*}

\end{document}
